\section{Trading Environment Modeling}

The proposed DRL framework is built upon custom \texttt{Gymnasium} environments 
that formulate the trading problem as a sequential decision-making task. 
At each timestep, the agent observes the current market state and portfolio 
information, and outputs continuous actions that determine asset allocation 
decisions. Two environments are implemented: a cryptocurrency trading environment
and a stock trading environment.

\subsection{Crypto Trading Environment}

The crypto trading environment models a multi-asset portfolio management problem, 
where the agent interacts with market observations over time and learns trading decisions
under realistic portfolio constraints. At each timestep $t$, the environment state is 
composed of two components: the portfolio state and the market state.

The portfolio state contains the available cash balance together with the current holdings
for all assets:
\[
s_t^{portfolio}
=
[C_t, H_1, H_2, \dots, H_N]
\]
where $N$ denotes the number of traded assets, $C_t$ is the available cash at timestep 
$t$, and $H_i$ represents the number of held shares (or units) of asset $i$.

The market state contains the current prices and technical indicators for all assets:
\[
s_t^{market}
=
[p_1, p_2, \dots, p_N,
\mathbf{tech}_1,
\mathbf{tech}_2,
\dots,
\mathbf{tech}_N]
\]
where $p_i$ is the current price of asset $i$, while
\[
\mathbf{tech}_i
=
[t_{i,1}, t_{i,2}, \dots, t_{i,F}]
\in \mathbb{R}^{F}
\]
is the feature vector containing the $F$ technical indicators associated with 
asset $i$. In practice, these indicators are flattened into a single vector and
concatenated with the asset prices to form the market representation. Therefore,
the market state combines both raw market information (prices) and engineered 
features (technical indicators).

The complete environment state at timestep $t$ is then defined as
\[
s_t
=
[s_t^{portfolio}, s_t^{market}]
\]
thus providing the agent with information about both its internal portfolio status and 
the external market conditions required for decision making.

The action space is continuous and bounded in $[-1,1]^N$, where each action component 
$a_i$ represents a trading signal for asset $i$. Negative values indicate selling 
pressure, while positive values indicate buying pressure.

For sell actions ($a_i < 0$), the environment liquidates a fraction of the currently 
held position proportional to the action magnitude:
\[
q_i^{sell} = \min(H_i,\,-a_iH_i)
\]
Therefore, $a_i=-1$ corresponds to full liquidation, while intermediate values produce 
partial sell orders.

For buy actions ($a_i > 0$), the action determines a desired allocation relative to the 
current portfolio value:
\[
V_t = C_t + \sum_{j=1}^{N} H_j p_j
\]

\[
q_i^{desired} = \frac{a_iV_t}{p_i}
\]

The final purchased quantity is constrained by available cash, transaction costs and 
maximum exposure limits:
\[
q_i^{buy} =
\min \left(
q_i^{desired},
\frac{C_t}{p_i(1+\kappa)},
\frac{\max(0,\lambda_iV_t-H_ip_i)}{p_i}
\right)
\]
where $\kappa$ denotes transaction costs and $\lambda_i$ defines the maximum portfolio 
allocation allowed for asset $i$. This prevents excessive concentration of capital 
into individual assets and enforces diversification constraints.

An alternative formulation considered during development was direct portfolio allocation, 
where actions lie in $[0,1]^N$ and represent target portfolio weights. However, this 
approach was discarded because it requires enforcing normalization constraints at each 
step and couples all asset decisions through the allocation vector. The adopted action 
representation allows independent asset-level decisions and scales more naturally to 
high-dimensional multi-asset portfolios.

The reward is defined as the logarithmic portfolio growth between consecutive timesteps:
\[
r_t =
\log\left(
\frac{V_{t+1}+\epsilon}
{V_t+\epsilon}
\right)
\]

Logarithmic returns were preferred over simple returns because they improve numerical 
stability, naturally account for compounding effects, and align the optimization process 
with long-term portfolio growth objectives.

Overall, the environment formulates crypto trading as a continuous-control MDP where the 
agent must maximize portfolio growth while accounting for transaction costs, capital 
constraints and position limits.

\subsection{Stock Trading Environment}

The stock trading environment extends the crypto formulation by incorporating additional 
market risk information and stricter trading constraints that better reflect traditional 
equity markets. Similar to the crypto environment, the state consists of a portfolio state 
and a market state, but the market representation is augmented with volatility and turbulence
signals.

The portfolio state remains identical:
\[
s_t^{portfolio}
=
[C_t,H_1,H_2,\dots,H_N]
\]
where $C_t$ denotes the available cash and $H_i$ the number of held shares of stock $i$.

The market state extends the previous formulation by including external market risk indicators:
\[
s_t^{market}
=
[p_1,\dots,p_N,
\mathbf{tech}_1,
\dots,
\mathbf{tech}_N,
v_t,
\tau_t]
\]
where $v_t$ denotes the market volatility index (VIX) at timestep $t$ and $\tau_t$ represents
the turbulence level of the market. The complete environment state is therefore given by
\[
s_t=
[s_t^{portfolio},s_t^{market}]
\]
allowing the agent to observe not only asset-level information but also global market risk
conditions.

The action space remains continuous and bounded in $[-1,1]^N$, where each component $a_i$
defines the trading signal for stock $i$. Sell actions reduce existing positions proportionally:
\[
q_i^{sell}
=
\min(H_i,-a_iH_i),
\qquad a_i<0
\]
while buy actions determine the desired position size relative to the current portfolio value
\[
V_t
=
C_t+\sum_{j=1}^{N}H_jp_j
\]

\[
q_i^{desired}
=
\left\lfloor
\frac{a_iV_t}{p_i}
\right\rfloor
\]

Unlike the crypto environment, stock holdings are constrained by an explicit maximum position
limit per asset. Therefore, the executed order becomes
\[
q_i^{buy}
=
\min\left(
q_i^{desired},
\left\lfloor
\frac{C_t}{p_i(1+\kappa)}
\right\rfloor,
M_i-H_i
\right)
\]
where $\kappa$ denotes transaction costs and $M_i$ the maximum allowable holdings for stock
$i$. This prevents unlimited accumulation of positions and enforces realistic exposure limits.

An additional risk-control mechanism is introduced through turbulence filtering. When market
turbulence exceeds a predefined threshold $\tau_{thr}$,
\[
\tau_t>\tau_{thr}
\]
all positive actions are suppressed according to
\[
a_t'=\min(a_t,0)
\]
forcing the agent to either maintain or reduce positions during periods of elevated market
stress.

The reward extends the logarithmic portfolio growth formulation by incorporating a volatility
penalty term:
\[
r_t
=
\log
\left(
\frac{V_{t+1}+\epsilon}
{V_t+\epsilon}
\right)
-
\lambda e_tv_t
\]
where
\[
e_t=
\frac{\sum_{i=1}^{N}H_ip_i}
{V_t+\epsilon}
\]
represents portfolio exposure, $v_t$ is the current market volatility signal (VIX), and
$\lambda$ controls the strength of the risk penalty. Consequently, portfolios with high
market exposure become increasingly penalized during volatile periods.

Overall, the stock environment formulates equity trading as a risk-aware continuous-control
MDP in which the agent must maximize long-term portfolio growth while accounting for trading
costs, position constraints, volatility exposure, and market turbulence.